% s-044-p044-m001; original PDF 44
\Aphorism{s-044-p044-m001}{30}
Es gehört zur Grossmuth des Edlen, alle Arten des Widerspruchs ohne Bitterkeit zu ertragen: den zwischen Schein und Wirklichkeit, Innen- und Aussenwelt, Tiefe und Oberfläche, Recht und Macht, Wollen und Können, eigener Bedeutung und fremder Schätzung — um nur einige Hauptfälle zu nennen. Umgekehrt bedarf der gallige rachsüchtige Charakter solcher Widersprüche als Gelegenheitsursachen, um ein aufgespeichertes Quantum von Ironie und Gehässigkeit loszuwerden. Man soll gegen alles Volk misstrauisch sein, das sich in der Nähe der Unlogik, Unästhetik, Unmoral zu schaffen macht; Früchte, die sich selbst so wenig Sonne gönnen, sind vielleicht noch unreif, vielleicht überhaupt missrathen, jedenfalls sauer und ungeniessbar.
\par

